% Generated draft: requires mathematical and source review.
\documentclass[11pt]{article}
\usepackage[margin=25mm]{geometry}
\usepackage{fontspec}
\setmainfont{Latin Modern Roman}
\usepackage{amsmath,amssymb,mathtools,unicode-math}
\setmathfont{Latin Modern Math}
\usepackage{enumitem,needspace,xurl,hyperref}
\setlength{\parindent}{0pt}
\setlength{\parskip}{5pt}
\setlength{\emergencystretch}{3em}
\newcommand{\N}{\mathbb{N}}
\newcommand{\Z}{\mathbb{Z}}
\newcommand{\Q}{\mathbb{Q}}
\newcommand{\R}{\mathbb{R}}
\newcommand{\C}{\mathbb{C}}
\newcommand{\F}{\mathbb{F}}
\newcommand{\A}{\mathbb{A}}
\newcommand{\PP}{\mathbb{P}}
\newcommand{\E}{\mathbb{E}}
\newcommand{\eps}{\varepsilon}
\newcommand{\dd}{\,\mathrm{d}}
\begin{document}
\section*{UnsolvedMath ID: 30001438}
\textbf{Source problem code:} catalogue entry
\subsection{Definitions and mathematical statement}
Let \(\Gamma\) be a countable discrete group, and let \(\beta^{(2)}_1(\Gamma)\) denote its first \(\ell^2\)-Betti number, defined as the von Neumann dimension over the group von Neumann algebra \(L\Gamma\) of its reduced first \(\ell^2\)-homology.

For a standard probability space \((K,\kappa)\), the associated Bernoulli action is the probability-measure-preserving action \(\Gamma\curvearrowright(X,\mu)=(K^\Gamma,\kappa^\Gamma)\) given by
\[
(g\cdot x)(h)=x(g^{-1}h)\qquad(g,h\in\Gamma).
\]
The source does not specify the base probability space, whether it must be nontrivial, or whether the question quantifies over all such bases. Accordingly, the formulation below retains the base as a parameter; it does not resolve this missing quantifier.

Under the standard interpretation, \(\mathrm{U}_{\mathrm{fin}}\) is the class of Polish groups isomorphic, as topological groups, to closed subgroups of \(\mathcal U(M)\), where \(M\) is a finite von Neumann algebra with separable predual and a faithful normal tracial state \(\tau\). The unitary group carries the topology induced by \(\|u-v\|_{2,\tau}=\tau((u-v)^*(u-v))^{1/2}\).

For a probability-measure-preserving action \(\Gamma\curvearrowright(X,\mu)\) and \(L\in\mathrm{U}_{\mathrm{fin}}\), a measurable cocycle is a measurable map \(c:\Gamma\times X\to L\) satisfying
\[
c(gh,x)=c(g,h\cdot x)c(h,x)
\]
for every \(g,h\in\Gamma\) and \(\mu\)-almost every \(x\). The action is \(\mathrm{U}_{\mathrm{fin}}\)-cocycle superrigid if, for every such \(L\) and \(c\), there exist a group homomorphism \(\rho:\Gamma\to L\) and a measurable map \(b:X\to L\) such that
\[
c(g,x)=b(g\cdot x)\rho(g)b(x)^{-1}
\]
for every \(g\in\Gamma\) and \(\mu\)-almost every \(x\).

\textbf{Source question, with the base qualification unresolved.} Does
\[
\beta^{(2)}_1(\Gamma)=0
\quad\Longrightarrow\quad
\Gamma\curvearrowright(K^\Gamma,\kappa^\Gamma)
\text{ is }\mathrm{U}_{\mathrm{fin}}\text{-cocycle superrigid}
\]
hold for countable groups \(\Gamma\), with the intended choice or quantification of \((K,\kappa)\) to be clarified? The supplied wording imposes no infinitude, nonamenability, or other additional group hypothesis; none is silently added here. Whether an additional restriction was intended requires source review.

The source singles out \(\Gamma=\mathbb Z\wr F_2\). Here \(F_2\) is the free group on two generators, and the usual restricted wreath product is
\[
\mathbb Z\wr F_2
=\left(\bigoplus_{h\in F_2}\mathbb Z\right)\rtimes F_2,
\qquad
(g\cdot f)(h)=f(g^{-1}h).
\]
Thus the singled-out question asks whether this group's Bernoulli action has the cocycle-superrigidity property defined above, subject to the same unresolved base qualification. The restricted-wreath-product and target-class definitions are standard interpretations of notation that the supplied record does not itself define.
\subsection{Short English statement}
The supplied question asks whether vanishing first \(\ell^2\)-Betti number forces every measurable cocycle of a Bernoulli action into a group in \(\mathrm{U}_{\mathrm{fin}}\) to be cohomologous to a homomorphism, especially for \(\mathbb Z\wr F_2\). The intended Bernoulli bases and any omitted group restrictions remain unresolved. The record labels the problem open; this draft does not independently verify that status.
\subsection{Sources}
\begin{itemize}
\item Local registry: \texttt{data/unsolvedmath/problems.json}, permanent UnsolvedMath ID \(30001438\), exact source problem code \texttt{catalogue entry}, title ``Cocycle Superrigidity for Bernoulli Actions.'' The supplied record belongs to the Oberwolfach Reports Open Problems collection and gives proposed year \(2010\); no individual proposer is identified.
\item Supplied primary citation: \emph{C*-Algebren} (2010), Oberwolfach Reports, DOI \texttt{10.4171/owr/2010/13}, pp.\ 667--732. \url{https://doi.org/10.4171/owr/2010/13}. The background describes these reports as records of mathematical workshops and their open-problem discussions. No precise page or speaker for this question is supplied.
\item Supplied literature source: ``Oberwolfach Report 13/2010.'' \url{https://doi.org/10.4171/owr/2010/13}.
\item Supplied literature source: ``Kerr--Li, entropy and orbit equivalence.'' \url{https://doi.org/10.1017/etds.2021.154}. This is the supplied reference label; no fuller bibliographic identification or supporting passage is provided.
\item Literature assessment dated \texttt{2026-08-21}: the supplied curation labels the problem open, states that the report identifies the Bernoulli \(\mathrm{U}_{\mathrm{fin}}\)-cocycle-superrigidity question for \(\mathbb Z\wr F_2\) as open, and says that later literature still describes nonamenable wreath-product cases as generally unresolved. This is supplied dated curation, not an independently verified assessment; the record expressly requests independent verification of time-sensitive status and cited claims.
\item The registry's statement-status label is \texttt{corrected\_verified}. Its accompanying verification note says only that TeX notation was enclosed in math delimiters without changing mathematical content. That note does not establish mathematical correctness or completeness of the formulation. No external source was visited in preparing this draft.
\end{itemize}
\end{document}
